\begin{pastexam}[label={px:ncku-106-ee-ds-03}]{成大電機所 106 學年度 複選題第 3 題}
（複選；每題 5 分，全對才給分）

Which of the following statement(s) is (are) true?
\begin{enumerate}[label=(\Alph*)]
\item A depth-first algorithm is implemented by queue.
\item The time complexity of the Prime algorithm is O(E logV) if it is implemented by an array.
\item There exists no back edge in a spanning tree generated by a breath-first algorithm.
\item A max heap can be implemented by a complete binary tree.
\end{enumerate}
\end{pastexam}

\begin{pastsoln}
答案：(C)(D)。

四個選項分屬資結書第~21~章（DFS、BFS）、第~24~章（Prim）、第~8~章（堆積），這裡放在第~21~章。 題目的
「Prime」是 Prim、「breath-first」是 breadth-first 的筆誤。

\textbf{(A) 錯。} 深度優先搜尋用的是\textbf{堆疊}：遞迴寫法用呼叫堆疊，迴圈寫法自己管一個堆疊（資結書
第~21.4.6~節）。 用佇列、先進先出的是廣度優先搜尋。

\textbf{(B) 錯。} Prim 演算法用陣列存每個頂點的 $\mathit{key}$、每次線性掃描找最小值的版本，是
$\Theta(V^2)$（資結書虛擬碼~24.5）：共 $V$ 次找最小值、每次 $\Theta(V)$，更新 $\mathit{key}$ 合計
$\Theta(E)$。 $O(E \log V)$ 是用二元堆積當優先佇列的版本。 兩者各有用處：稠密圖 $E \approx V^2$ 時陣列版
$\Theta(V^2)$ 反而比 $O(V^2 \log V)$ 好。

\textbf{(C) 對。} 生成樹裡只有樹邊，本來就沒有非樹邊。 再從「BFS 會不會產生後向邊」看：對無向圖做 BFS，
一條非樹邊 $(u, v)$ 兩端的層數最多差 $1$（資結書第~21.3~節）。 若 $u$ 是 $v$ 的祖先、層數差 $1$，$v$ 就是 $u$
的子節點，這條邊就是樹邊本身（沒有重邊時），所以非樹邊不會連接祖先和子孫，只能是交叉邊，沒有後向邊。
（有向圖的 BFS 可以有後向邊，例如有向環 $1 \to 2 \to 1$；題目講的生成樹是無向圖的概念。）

\textbf{(D) 對。} 二元堆積就是一棵滿足堆積性質（每個節點不小於子節點）的完整二元樹，通常用陣列存
（資結書第~8.1~節）。
\end{pastsoln}
